% Imporditud: tools/import_eksperimendid.py
% Allikas: Eesti füüsikaolümpiaadi eksperimendi ülesannete
%   kogu (Rael Kalda, 2023), ülesanne Ü132, lahendus L132.
\setAuthor{EFO žürii}
\setRound{piirkonnavoor}
\setYear{2006}
\setNumber{G E2}
\setDifficulty{4}
\setTopic{dünaamika}
\setVahendid{millimeeterskaalaga tikutoosi kest, paberileht, tikk}
\setPunktid{14}
\setAllikas{Eksperimendi kogu Ü132, lahendus lk 102}
\setLahendusseis{toores}

\prob{Tikutoos}
Määrake tikutoosi kesta otsa ja paberi vaheline hõõrdetegur. Tikutoosi mõõtmed on 50 $\times$ 36 $\times$ 15 mm.

\hint

\solu
Tikutoosi paberil tikuga lükates mõjub sellele jõud F = $\mu$mg (1 p). Teatud punktist altpoolt lükates tikutoos libiseb ning ülevalt poolt lükates kukub tikutoos ümber (1 p). Libisemise tingimus:

mgd Fl $\leq$ ,

kus l on kõrgus, millelt lükatakse, ja d tikutoosi paksus, m mass ja g raskuskiirendus. Ehk siis jõumoment, mis pöörab tikutoosi algasendisse tagasi, on suurem kui tikutoosi ümber lükkav jõumoment (5 p). Kasutades suurimat kõrgust, millelt lükates tikutoos veel libiseb, ehk

d (1 p) $\Rightarrow$ d (0,5 p) $\Rightarrow$ d (0,5 p) F l = mg $\mu$mgl = mg $\mu$= . 2 2 2l

Maksimaalse kõrguse määramine katseliselt, millelt tikutoos veel libiseb (3 p), sest see ei ole sugugi nii lihtne. Korduskatsed ja mingisugunegi veahinnang, näiteks: mõõtmistulemuste aritmeetiline keskmine (2 p).
\probend
